\documentclass{beamer}
\usepackage{beamerthemesplit}
\usepackage{wrapfig}
\usetheme{SPbGU}
\usepackage{pdfpages}
\usepackage{amsmath}
\usepackage{cmap} 
\usepackage[T2A]{fontenc} 
\usepackage[utf8]{inputenc}
\usepackage[english,russian]{babel}
\usepackage{indentfirst}
\usepackage{amsmath}
\usepackage{tikz}
\usepackage{multirow}
\usepackage[noend]{algpseudocode}
\usepackage{algorithm}
\usepackage{algorithmicx}
\usepackage{listings}

\usepackage{epstopdf}
\usepackage{forest}
\usetikzlibrary{shapes,arrows}
\usepackage{fancyvrb}
\newtheorem{rutheorem}{Theorem}
\beamertemplatenavigationsymbolsempty

\lstdefinelanguage{ocaml}{
keywords={fresh, let, begin, end, in, match, type, and, fun, function, try, with, when, class, 
object, method, of, rec, repeat, until, while, not, do, done, as, val, inherit, 
new, module, sig, deriving, datatype, struct, if, then, else, open, private, virtual, ostap},
sensitive=true,
basicstyle=\small,
commentstyle=\small\itshape\ttfamily,
keywordstyle=\ttfamily\underbar,
identifierstyle=\ttfamily,
basewidth={0.5em,0.5em},
columns=fixed,
fontadjust=true,
literate={->}{{$\;\;\to\;\;$}}1,
morecomment=[s]{(*}{*)}
}

\lstset{
basicstyle=\small,
identifierstyle=\ttfamily,
keywordstyle=\bfseries,
commentstyle=\scriptsize\rmfamily,
basewidth={0.5em,0.5em},
fontadjust=true,
%escapechar=~,
language=ocaml,
mathescape=true
}

\title[]{Ostap: синтаксическое расширение OCaml для создания парсер-комбинаторов с поддержкой левой рекурсии}
\institute[СПбГУ]{
Санкт-Петербургский государственный университет\\
Лаборатория языковых инструментов JetBrains }

\author[Екатерина Вербицкая]{Екатерина Вербицкая}

\date{05 апреля 2017}

\definecolor{orange}{RGB}{179,36,31}

\begin{document}
{

\begin{frame}
  \begin{center}
  {\includegraphics[width=1cm]{SPbGU_Logo.png}}
  \end{center}
  \titlepage
\end{frame}
}

\begin{frame}[fragile]
  \transwipe[direction=90]
  \frametitle{Синтаксический анализ}
  Сопоставление последовательности лексем с грамматикой языка
  
  \begin{columns}
    \begin{column}{8cm}
  \begin{minipage}[t]{8cm}

\begin{center}
 $1 + 2 * 3$
\end{center}

\medskip


$$
\begin{array}{crcl}
&e & : & e + t \ | \ t \\
&t & : & t * p \ | \ p \\
&p & : & 0 \ | \ 1 \ | \ 2 \ | \dots | \ 9
\end{array}
$$
\end{minipage}
\end{column}
\begin{column}{5cm}
\begin{forest}
  [e[e[t[p[$1$]]]][$+$][t[t[p[$2$]]][$*$][p[$3$]]]]
\end{forest}
\end{column}
\end{columns}
\end{frame}

\begin{frame}
  \transwipe[direction=90]
  \frametitle{Парсер-комбинаторы}  
  \begin{itemize}
  \item Подход к реализации синтаксического анализа в парадигме функционального программирования
  \item Реализуют нисходящий разбор
  \item Синтаксический анализатор: функция высшего порядка
  \item Позволяет считать семантику ``на лету'', без явного построения деревьев разбора
  \item Анализатор произвольной сложности можно получить путем комбинирования нескольких простых базовых парсеров
  \item Позволяют разбирать некоторые контекстно-зависимые языки
  \end{itemize}
\end{frame}

\begin{frame}[containsverbatim]
  \transwipe[direction=90]
  \frametitle{Тип парсеров}
  
\begin{lstlisting}[frame=single]  
type $\alpha$ tag = Parsed of $\alpha$ | Failed

type ($\sigma$, $\rho$) result = ($\rho$ * $\sigma$) tag

and  ($\sigma$, $\rho$) parse  = $\sigma$ ->   ($\sigma$, $\rho$) result
\end{lstlisting}
\end{frame}


\begin{frame}[fragile]
  \transwipe[direction=90]
  \frametitle{Базовые парсер-комбинаторы}  
\begin{lstlisting}[frame=single]  
let empty s = Parsed ((), s)

let fail  s = Failed 

let return x s = Parsed (x, s)
\end{lstlisting}
\end{frame}

\begin{frame}[fragile]
  \transwipe[direction=90]
  \frametitle{Комбинатор последовательности}  
\begin{lstlisting}[frame=single]  
let seq x y s =
  match x s with
  | Parsed (b, s') -> y b s'
  | Failed -> Failed

let (|>) = seq
\end{lstlisting}
\end{frame}


\begin{frame}[fragile]
  \transwipe[direction=90]
  \frametitle{Комбинатор альтернативы}  
\begin{lstlisting}[frame=single]  
let alt x y s =
  match x s with
  | Failed x -> y s
  | x -> x

let (<|>) = alt
\end{lstlisting}
\end{frame}


\begin{frame}[fragile]
  \transwipe[direction=90]
  \frametitle{Пример парсера для языка $\{ a^n \ | \ n \geq 0 \}$}  
\begin{lstlisting}[frame=single]  
(* A : epsilon | 'a' A *)
let rec a s = empty <|> (terminal 'a' |> fun _ -> a) s
\end{lstlisting}
\end{frame}


\begin{frame}[fragile]
  \transwipe[direction=90]
  \frametitle{Семантические действия}  
\begin{lstlisting}[frame=single]  
let map f p s = 
  match p s with
  | Parsed (b, s') -> Parsed (f b, s')
  | x -> x
\end{lstlisting}
\end{frame}

\begin{frame}[fragile]
  \transwipe[direction=90]
  \frametitle{Вспомогательные парсер-комбинаторы}  
\begin{lstlisting}[frame=single]  
let opt p = map (fun x -> Some x) p <|> return None
    
let rec many p =
  p |> (fun h -> map (fun t -> h :: t) (many p) <|> return [])
    
\end{lstlisting}

\begin{lstlisting}[frame=single]  
let a = many (terminal 'a')
\end{lstlisting}
\end{frame}

\begin{frame}[fragile]
  \transwipe[direction=90]
  \frametitle{Синтаксическое расширение}  
\begin{lstlisting}[frame=single]  
ostap (
  e : x:e "+" y:t {Add (x, y)} 
    | t; 
  t : x:t "*" y:p {Mul (x, y)} 
    | p;
  p : (* parse integer *)
)
\end{lstlisting}
\end{frame}

\begin{frame}[fragile]
  \transwipe[direction=90]
  \frametitle{Парсеры высшего порядка: переиспользование кода}  
\begin{lstlisting}[frame=single]  
ostap (
  e : x:e "+" y:t {Add (x, y)} 
    | t; 
  t : x:t "*" y:p {Mul (x, y)} 
    | p;
  p : (* parse integer *)
)
\end{lstlisting}
  
\begin{lstlisting}[frame=single]  
ostap (
  e : x:e "+" y:t {Add (x, y)} 
    | t; 
  t : x:t "*" y:p {Mul (x, y)} 
    | p;
  p : (* parse double *)
)
\end{lstlisting}	  
  
\end{frame}


\begin{frame}[fragile]
  \transwipe[direction=90]
  \frametitle{Парсеры высшего порядка}  
\begin{lstlisting}[frame=single]  
ostap (
  e[p] : x:e[p] "+" y:t[p] {Add (x, y)} 
       | t[p]; 
  t[p] : x:t[p] "*" y:p    {Mul (x, y)} 
       | p;
)
\end{lstlisting}
  
\begin{lstlisting}[frame=single]  
ostap (
  e_integer: e[(* parse integer *)]
)
\end{lstlisting}
  
\begin{lstlisting}[frame=single] 
ostap (
  e_double : e[(* parse double  *)]
)
\end{lstlisting}
\end{frame}


\begin{frame}[fragile]
  \transwipe[direction=90]
  \frametitle{Парсер-комбинаторы и левая рекурсия}  
Являясь реализацией нисходящего анализа, парсер-комбинаторы не способны обрабатывать леворекурсивные правила

\begin{lstlisting}[frame=single]  
ostap (
  e[p] : e[p] "+" t[p] (* to apply e, *)
       | t[p];         (* one needs to apply e... *)
  t[p] : t[p] "*" p    
       | p;
)
\end{lstlisting}

Удаление левой рекурсии значительно усложняет спецификацию языка и ухудшает композициональность анализаторов
\end{frame}


\begin{frame}[fragile]
  \transwipe[direction=90]
  \frametitle{Поддержка леворекурсивных анализаторов}
  \begin{itemize} % СНАЧАЛА СУТЬ, ПОТОМ КРЕДИТ, ПРИ ЭТОМ БЕЗ НАЗВАНИЯ, ТОЛЬКО ГОД
    \item Ограничение количества леворекурсивных вызовов длиной непрочитанной строки
    \begin{itemize}
      \item Frost R. A., Hafiz R., Callaghan P. Parser combinators for ambiguous leftrecursive grammars
    \end{itemize}
    \item Использование мемоизации для обеспечения завершаемости
    \begin{itemize}
      \item Warth A., Douglass J. R., Millstein T. D. Packrat parsers can support left recursion
    \end{itemize}

    \item Требуют, чтобы парсер был первого порядка
  \end{itemize}

  \begin{itemize}
    \item Использование техники CPS для обеспечения завершаемости
    \begin{itemize}
      \item Izmaylova	A., Afroozeh A., Tijs van der Storm. 
Practical, general parser combinators
    \end{itemize}

    \item Фиксируют конкретную семантику
  \end{itemize}  

\end{frame}

%\begin{frame}[fragile]
%  \transwipe[direction=90]
%  \frametitle{Полиморфность относительно типа входного потока}
%\begin{itemize}
%  \item Позволяет инкорпорировать в себя сканер (лексер)
%  \item Повышает композиционность еще сильнее
%  \item Позволяет реализовывать в абстракции входного потока другие полезные фичи
%\end{itemize}

%\end{frame}


\begin{frame}
  \transwipe[direction=90]
  \frametitle{Поддержка левой рекурсии в PEG}
  \begin{itemize}
    \item Medeiros S., Mascarenhas F., Ierusalimschy R. Left Recursion in Parsing Expression Grammars 
    \item Динамический поиск наилучшего количества леворекурсивных вызовов
    \item Использует мемоизацию
    \item Поддерживает явную, неявную, взаимную рекурсию
    \item Требуют, чтобы парсер был первого порядка
  \end{itemize}
\end{frame}

\begin{frame}[fragile]
  \transwipe[direction=90]
  \frametitle{Поддержка левой рекурсии в Ostap}
  \begin{itemize} 
    \item Используется идея Medeiros et al
    \item Специальный комбинатор \verb!fix! для поддержки леворекурсивных парсеров высшего порядка
  \end{itemize}
\begin{lstlisting}[frame=single]  
ostap (
  e[p] : e[p] "+" e[p] | p
)
\end{lstlisting}

\begin{lstlisting}[frame=single]  
ostap (
  let e p s = fix (fun e -> ostap (e "+" e | p)) s
)
\end{lstlisting}
\end{frame}


\begin{frame}[fragile]
  \transwipe[direction=90]
  \frametitle{Заключение}
  \begin{itemize} 
    \item Представлена библиотека парсер-комбинаторов и синтаксическое расширение для языка OCaml
    \item Реализована поддержка леворекурсивных спецификаций синтаксических анализаторов
    \begin{itemize}
      \item Позволяет использование парсеров высшего порядка
      \item Сложность растет экспоненциально в зависимости от глубины вложенности рекурсии
    \end{itemize}
  \end{itemize}
\end{frame}


\end{document}